The two geometric pictures now have a common algebraic carrier.  A one-hole
arithmetic section is represented by a $2\times2$ matrix; chronological
composition is matrix multiplication, point paths are matrix orbits, ripple
pencils are inverse images under the same products, poles are inverse images
of infinity, and fixed points are eigenlines.

\subsection{Elementary arithmetic matrices}

Let $K$ be a field.  In homogeneous coordinates, write
\[
  z=[z:1],
  \qquad
  \infty=[1:0],
\]
and let
\[
  M=\begin{pmatrix}A&B\\C&D\end{pmatrix}
\]
act by
\[
  M\cdot z=\frac{Az+B}{Cz+D}.
\]
Nonzero scalar multiples represent the same element of $\PGL_2(K)$.

The nondegenerate elementary sections have the following representatives.

\begin{center}
\renewcommand{\arraystretch}{1.4}
\begin{tabular}{lll}
\toprule
one-hole section & affine formula & matrix \\
\midrule
$C^{(1)}_{+,c}=C^{(2)}_{+,c}$
& $z+c$
& $\begin{pmatrix}1&c\\0&1\end{pmatrix}$ \\
$C^{(1)}_{-,c}$
& $z-c$
& $\begin{pmatrix}1&-c\\0&1\end{pmatrix}$ \\
$C^{(2)}_{-,c}$
& $c-z$
& $\begin{pmatrix}-1&c\\0&1\end{pmatrix}$ \\
$C^{(1)}_{\times,c}=C^{(2)}_{\times,c}$
& $cz$
& $\begin{pmatrix}c&0\\0&1\end{pmatrix}$ \\
$C^{(1)}_{/,c}$
& $z/c$
& $\begin{pmatrix}1&0\\0&c\end{pmatrix}$ \\
$C^{(2)}_{/,c}$
& $c/z$
& $\begin{pmatrix}0&c\\1&0\end{pmatrix}$ \\
\bottomrule
\end{tabular}
\end{center}

Here $c\ne0$ whenever invertibility requires it.  The table has two immediate
consequences.

\begin{enumerate}
  \item Every pure slot-$(1)$ section is affine and fixes $\infty$.
  \item The projectively new primitive is slot-$(2)$ division.  Up to a
  nonzero scaling, it is inversion.
\end{enumerate}

Ordinary arithmetic and projective continuation remain separate.  A matrix in
$\PGL_2(K)$ acts everywhere on $\PP^1(K)$, while the corresponding ordinary
word may be undefined at an intermediate denominator zero.

\subsection{Chronological products}

Let $F_1,\ldots,F_n$ be nondegenerate sections with matrix representatives
$M_1,\ldots,M_n$, and suppose $F_1$ is applied first.  Define
\begin{equation}\label{eq:p0-prefix-matrices}
  P_k=M_kM_{k-1}\cdots M_1,
  \qquad
  G_k=F_k\circ\cdots\circ F_1.
\end{equation}
Then $P_k$ represents $G_k$.

\begin{theorem}[One matrix, four readings]
\label{thm:p0-four-matrix-readings}
Write
\[
  P_k=\begin{pmatrix}A_k&B_k\\C_k&D_k\end{pmatrix}.
\]
The same prefix matrix determines:

\begin{enumerate}
  \item the point-path state
  \[
    z_k=G_k(z_0)=\frac{A_kz_0+B_k}{C_kz_0+D_k};
  \]

  \item the projective pole
  \[
    p_k=G_k^{-1}(\infty)
      =-D_k/C_k
  \]
  when $C_k\ne0$, and $p_k=\infty$ when $C_k=0$;

  \item the fixed-point equation
  \[
    C_kz^2+(D_k-A_k)z-B_k=0;
  \]

  \item over $\R$, when a lift has positive determinant, the ripple pencil
  $G_k^{-1}(\mathcal H_t)$ from
  Theorem~\ref{thm:p0-ripple-pencil-formula}.
\end{enumerate}
\end{theorem}

\begin{proof}
The first formula is the matrix action.  The pole solves
$C_kz+D_k=0$.  The fixed-point equation comes from
$z(C_kz+D_k)=A_kz+B_k$.  The final statement is the definition of the ripple
pencil together with formula \eqref{eq:p0-imaginary-mobius}.
\end{proof}

Thus \emph{path} and \emph{ripple} are not two unrelated representations.  They
are the covariant and contravariant geometries of one chronological matrix
product.

\begin{figure}[t]
\centering
\begin{tikzpicture}[
  box/.style={draw=black!60,rounded corners=3pt,fill=blue!5,
              minimum width=31mm,minimum height=9mm,align=center,font=\small},
  arr/.style={-{Latex[length=2mm]},thick,black!65},
  every node/.style={font=\small}
]
  \node[box,fill=violet!7] (word) at (0,2.0)
    {one-hole word\\$F_1,\ldots,F_k$};
  \node[box,fill=orange!8] (matrix) at (0,0.65)
    {prefix matrix\\$P_k=M_k\cdots M_1$};
  \node[box] (path) at (-4.2,-1.1) {point path\\$P_k\cdot z_0$};
  \node[box] (pole) at (-1.4,-1.1) {pole\\$P_k^{-1}(\infty)$};
  \node[box] (fixed) at (1.4,-1.1) {fixed points\\eigenlines of $P_k$};
  \node[box] (ripple) at (4.2,-1.1) {ripple pencil\\$P_k^{-1}(\mathcal H_t)$};
  \draw[arr] (word)--(matrix);
  \foreach \n in {path,pole,fixed,ripple}
    {\draw[arr] (matrix)--(\n);}
\end{tikzpicture}
\caption{Four geometric readings of the same chronological matrix product.}
\label{fig:p0-projective-four-readings}
\end{figure}

\subsection{The affine chart as the stabilizer of infinity}

Let
\[
  B_\infty=\Stab_{\PGL_2(K)}(\infty).
\]
A matrix fixes infinity exactly when $C=0$.  After projective rescaling, every
such element has a unique representative
\[
  \begin{pmatrix}a&b\\0&1\end{pmatrix},
  \qquad a\in K^\times,\quad b\in K,
\]
and acts as $z\mapsto az+b$.  Hence
\begin{equation}\label{eq:p0-affine-stabilizer}
  B_\infty\cong\Aff(1,K).
\end{equation}

This identifies the path geometry of Section~\ref{sec:p0-paths} as a
pole-at-infinity gauge.  In that gauge:

\begin{itemize}
  \item the denominator is constant;
  \item the ripple pencil is a parallel-line pencil;
  \item fixed-point recursion is affine;
  \item the lower-left matrix entry stays zero.
\end{itemize}

Choosing a finite pole $q$ amounts to conjugating this affine subgroup by a
projective map sending $q$ to infinity.  The coordinate
$w_q=1/(z-q)$ from \eqref{eq:p0-pole-coordinate} is the simplest such gauge.

\subsection{Generation of the projective group}

Introduce
\[
  T_s(z)=z+s,
  \qquad
  D_q(z)=qz\quad(q\ne0),
  \qquad
  J(z)=-\frac1z.
\]
They arise respectively from addition, multiplication, and second-slot
division.

\begin{theorem}[Elementary projective generation]
\label{thm:p0-pgl2-generation}
For every field $K$, translations $T_s$, nonzero scalings $D_q$, and inversion
$J$ generate $\PGL_2(K)$.  Equivalently, the projective evaluations of
nondegenerate left- and right-slot arithmetic sections have image
$\PGL_2(K)$.
\end{theorem}

\begin{proof}
Let
\[
  F(z)=\frac{Az+B}{Cz+D},
  \qquad AD-BC\ne0.
\]
If $C=0$, then $A,D\ne0$ and
\[
  F=T_{B/D}\circ D_{A/D}.
\]
If $C\ne0$, direct calculation gives
\begin{equation}\label{eq:p0-projective-factorization}
  F
  =T_{A/C}
   \circ D_{(AD-BC)/C^2}
   \circ J
   \circ T_{D/C}.
\end{equation}
All scale parameters are nonzero, so every factor is available.  These two
cases exhaust $\PGL_2(K)$.
\end{proof}

The proof also yields the rank-one Bruhat decomposition
\begin{equation}\label{eq:p0-bruhat}
  \PGL_2(K)=B_\infty\sqcup B_\infty J B_\infty.
\end{equation}
The first cell consists of operators whose pole is infinity.  The second
consists of operators with a finite pole in the chosen affine chart.  Thus the
transition from a point-path grid to a curved ripple pencil is the elementary
geometric face of crossing from the affine cell to the projective cell.

\subsection{Left and right representations in one projective plane}

We can now answer the algebraic comparison of the two combs precisely.

\begin{theorem}[Projective unification of the two combs]
\label{thm:p0-comb-projective-unification}
The left- and right-expanded combs have the same abstract chain of internal
dependencies and are exchanged by the opposite-operation involution of
Proposition~\ref{prop:p0-comb-opposition}.  Their one-hole operator images have
the following relation:

\begin{enumerate}
  \item the pure left-expanded nondegenerate language lies in
  $B_\infty\cong\Aff(1,K)$;

  \item the pure right-expanded language already contains translations,
  nonzero scalings, reflections, and inversion, and therefore generates
  $\PGL_2(K)$;

  \item after embedding both languages into $\PGL_2(K)$, their point paths,
  poles, fixed points, and ripple pencils are all read from the same matrix
  representation.
\end{enumerate}
\end{theorem}

\begin{proof}
The dependency and opposition statements were proved in
Section~\ref{sec:p0-two-combs}.  The elementary table shows that every
slot-$(1)$ generator fixes infinity.  The pure slot-$(2)$ language contains
$T_s$ because addition is commutative, contains $D_q$ because multiplication
is commutative, and contains $J$ through second-slot division.  Apply
Theorem~\ref{thm:p0-pgl2-generation}.  The last assertion is
Theorem~\ref{thm:p0-four-matrix-readings}.
\end{proof}

The theorem does not say that a left-expanded word and a right-expanded word
with the same constants have the same operator.  It says that their apparently
different elementary geometries live in one projective action and can be
compared there without erasing operand-slot information.
